
\documentclass[a4paper,11pt,twocolumn]{article}

\usepackage[utf8]{inputenc}
\usepackage[T1]{fontenc}
\usepackage[dvips]{graphicx}
%\usepackage{framed}
\usepackage{latexsym}
\usepackage{amsmath}
\usepackage{amssymb}
\usepackage{listings}
\usepackage{courier}
\usepackage{caption}
\usepackage[normalem]{ulem}
\usepackage{pgfplots}
\usepackage{booktabs,colortbl}
\usepackage{hyperref}
\usepackage{url}
\usepackage{etex}
\usepackage{pstricks-add}
\usepackage{paralist}
\usepackage[nounderscore]{syntax}
\usepackage{tikz}
\usetikzlibrary{matrix,shapes,shapes.multipart,backgrounds,calc,fit}

%\definecolor{shadecolor}{gray}{0.9}
 
\setlength{\fboxsep}{1.5pt}
\newcommand{\highl}{\colorbox{gray!30}}

\lstloadlanguages{C}

\def\lstsmallmath{\leavevmode\ifmmode \scriptstyle \else  \fi}
\def\lstsmallmathend{\leavevmode\ifmmode  \else  \fi}
\lstset {
  escapechar=§,
  %frame=shadowbox,
  %rulesepcolor=\color{black},
  showspaces=false,showtabs=false,tabsize=2,
  numberstyle=\tiny,numbers=left,
  stringstyle=\color{sh_string},
  keywordstyle = \color{sh_keyword}\bfseries,
  commentstyle=\color{sh_comment}\itshape,
  captionpos=b,
  xleftmargin=0.7cm,
  aboveskip=0.7cm,
  belowskip=0.5cm,
  escapebegin={\lstsmallmath}, escapeend={\lstsmallmathend},
  emph={let,and,match,function,with,in,assert,true,false,None,Some,fun,
    int,char,void,return,size\_t,unsigned,long,struct,if,then,else,raise,type,
    method,module,end,rec,ref,val,list,for,while,break,const,short},
  emphstyle=\textbf
}


\captionsetup[lstlisting]{
  singlelinecheck=false,
  margin=40pt,
  font={bf,footnotesize}
}



\begin{document}
 
\title{Code generation for memory monitoring}
\author{...}
\date{\today}
 
\maketitle




\section{Introduction}


Memory accesses from a pointer or an array are not safe in C, thus trying to
access an element of an array out of valid range or an invalid pointer may cause
a segmentation fault during execution.

C compilers such as \textsc{Gcc} do not provide detection mechanisms for this
kind of errors. We propose an automatic code instrumentation, so that the
generated code will perform memory monitoring and will be able to check
memory properties during execution (runtime checking).
We consider as readable and writable memory: global variables, dynamically
allocated memory, and formal parameters and local variables of each function. We
decided to use the term ``block'' to nominate these four cases.

\textsc{E-Acsl} is an executable subset of \textsc{Acsl} \cite{ACSL}, a formal
specification language for C using source code annotations to express
properties. \textsc{Acsl}-annotated C programs can be dealt with
\textsc{Frama-C} \cite{Frama-C}, a framework for modular analysis of C.
\textsc{Frama-C} provides a plug-in generating executable C code from
\textsc{E-Acsl} annotations, used by the implementation discussed in this paper.
Memory blocks (as defined above) held properties such as their base address and
their size. \textsc{E-Acsl} provides five annotations to retrieve useful
information about blocks:

\begin{description}
\item[$\backslash$base\_addr(p)] \hfill \\
  returns the base address of the block containing the pointer $p$
\item[$\backslash$block\_length(p)] \hfill \\
  returns the size (in bytes) of the block containing the pointer $p$
\item[$\backslash$offset(p)] \hfill \\
  returns the offset between $p$ and $\backslash$base\_addr($p$)
\item[$\backslash$valid(s)] \hfill \\
  whether reading and writing $*s$ is safe
\item[$\backslash$initialized(s)] \hfill \\
  whether the variable stored at the address $s$ has been initialized
\end{description}

Our contribution allowed the \textsc{E-Acsl} plug-in to generate executable C
code from these five annotations, thus checking these five memory properties
during execution.



\section{Stored information}

For each block, we need to store the following information: the base address
(address of the first element within the block), the size (in bytes), the
validity status (whether reading and writing the block is safe) and the
initialization status for each byte of the block.
Figure~\ref{fig:block_descr} illustrates the data structure (we call it
{\em block descriptor}) chosen to store this information.


\begin{figure}[h]
  \begin{lstlisting}
  struct _block {
    char * ptr;
    size_t size;
    int valid;
    unsigned char * init_ptr;
    unsigned long init_cpt;
  };
  \end{lstlisting}
  \caption{Block descriptor}
  \label{fig:block_descr}
\end{figure}


$init\_ptr$ is an array of booleans that is dynamically allocated only if
needed. If it has been allocated, it contains a boolean for each byte of the
block: the $n^{th}$ boolean indicates whether the $n^{th}$ byte has been
initialized.

$init\_cpt$ counts the number of initialized bytes within the block. If
$init\_cpt = 0$ (none) or $init\_cpt = size$ (all) then $init\_ptr$ is freed.
So, when none or all of the bytes have been initialized (the most common cases),
the memory space needed for the block descriptor itself is reduced. This
consistency is maintained when adding/removing an element.


\section{Global architecture}

\newcommand{\colB}{0}
\newcommand{\rowA}{7}
\newcommand{\rowB}{5}
\newcommand{\rowC}{3}
\newcommand{\rowD}{1}
\begin{figure}[h]
  \begin{center}
    \tikzstyle{service}=[rounded corners,rectangle,inner sep=2mm,fill=gray!40]
    \scalebox{0.85}{
      \begin{tikzpicture}
        \node[service,rectangle,inner sep=4mm] (eacsl) at (\colB{},\rowA{}) {
          E-ACSL plug-in};
        \node [inner sep=4mm,draw](lib) at (\colB{},\rowB{})
              {Monitoring functions};
        \node[inner sep=4mm,draw](sdda) at (\colB{},\rowC{}){Common interface};
        \node [matrix of nodes,rectangle,draw,column sep=4mm,inner sep=2.5mm]
        (sddc) at (\colB{},\rowD{}) {
          \node[service](p) at (0,0) {{\em Bittree}}; \pgfmatrixnextcell
          \node at (0,0) {or}; \pgfmatrixnextcell
          \node[service] at (0,0) { Linked list }; \pgfmatrixnextcell
          \node at (0,0) {\dots};\\
        };
        \path [<->,thick] (eacsl) edge (lib);
        \path [<->,thick] (lib) edge (sdda);
        \path [<->,thick] (sdda) edge (sddc);
      \end{tikzpicture}
    }
  \end{center}
  \caption{Global architecture}
  \label{fig:lib}
\end{figure}


Interactions between the components of our solution are displayed by
Figure~\ref{fig:lib}. A data structure (see
Section~\ref{section:bittree}) stores the block descriptors. The library
functions (see Section~\ref{section:lib}) add, remove or modify the block
descriptors. The \textsc{E-Acsl} plug-in then perform an instrumentation (see
Section~\ref{section:instru}) to generate executable C code. The generated code
uses the monitoring functions of this library.


\section{Concrete data structure used: {\em bittree}}
\label{section:bittree}

In this section we discuss the choice of the concrete data structure used to
store the block descriptors.


\subsection{{\em Patricia tries}}


We need a data structure with a good time and space complexity, indeed we may
have to often add or remove a block descriptor in the structure. The structure
has to be sorted: we want to access to a block descriptor by its base address,
and also to its predecessor and successor. Thus, hash-tables will not fit.
We also unconsidered the linked lists, due to the linear worst-case complexity.
The (unbalanced) binary search trees provide a linear worst-case complexity too
when the base address of inserted elements are monotonically increasing, and
this may be quite common. Self-balanced binary search trees are dismissed
because of the numerous add/remove operations implying numerous costly balancing
operations.

We choose to use the {\em Patricia tries} \cite{Patricia-tries} structure, which
is efficient even if the tree is unbalanced. The prefix used by a node (not a
leaf) is the greatest common prefix (on 32 or 64 bits) of its two children.
The block descriptors are held on leaves. Other nodes just do the routing from
the root to a block descriptor. Patricia tries are usually used on
strings and characters, so we named ``bittree'' the structure used in this
article. For example on 8-bit addresses, Figure~\ref{fig:bittree} shows
a bittree  storing three block descriptors (identified by their addresses:
${\tt 0010\,0111}$, ${\tt 0010\,1001}$ and ${\tt 0010\,1101})$. The greatest
common prefix of ${\tt 0010\,1001}$ and ${\tt 0010\,1101}$ is ${\tt 0010\,1***}$
and the greatest common prefix of ${\tt 0010\,0111}$ and ${\tt 0010\,1***}$ is
${\tt 0010\,****}$. The $*$ means that this bit is meaningless.
Figure~\ref{fig:insertion} shows that a new node (with a new prefix) is added
when a block descriptor is inserted.

Figure~\ref{fig:deletion} shows an example of deletion of a 8-bit block
descriptor. Patricia tries being compact prefix trees, a node having an
only child is deleted. So ${\tt 0010\,1***}$ becomes useless and is replaced by
its child ${\tt 0010\,1001}$. Deleting useless nodes, or only storing useful
ones for routing, keeps the tree exploration efficient. A 32-bit (respectively
64-bit) bittree has a worst-case depth of 33 (respectively 65) nodes.


\begin{figure}[h]
  \begin{center}
    \tikz[grow=right,level 2/.style={sibling distance=1em},level distance=3cm]
    \node {0010 ****}
    child { node {0010 0111} }
    child { node {0010 1***}
      child { node {0010 1001} }
      child { node {0010 1101} }
    };
    \caption{Example of bittree}
    \label{fig:bittree}
  \end{center}
\end{figure}

\begin{figure}[h]
  \begin{center}
    \tikz[grow=right,level 2/.style={sibling distance=1em},level distance=3cm]
    \node {0010 ****}
    child { node {0010 011*}
      child { node {0010 0110} }
      child { node {0010 0111} }
    }
    child { node {0010 1***}
      child { node {0010 1001} }
      child { node {0010 1101} }
    };
    \caption{Bittree after adding ${\tt 0010\,0110}$ into the bittree of
      Fig.~\ref{fig:bittree}}
    \label{fig:insertion}
  \end{center}
\end{figure}

\begin{figure}[h]
  \begin{center}
    \tikz[grow=right,level 2/.style={sibling distance=1em},level distance=3cm]
    \node {0010 ****}
    child { node {0010 0111} }
    child { node {0010 1001} };
    \caption{Bittree after deleting ${\tt 0010\,1101}$ from the bittree of
      Fig.~\ref{fig:bittree}}
    \label{fig:deletion}
  \end{center}
\end{figure}


\subsection{Greatest common prefix computation}
\label{section:gcpc}

The greatest common prefix of $A$ and $B$ can be naively computed by:
$X = \lnot (A \oplus B)$ (where $\oplus$ is the XOR operator) to keep bits in
common, then each bit of $X$ on the right side of a 0 is set to 0. The obtained
mask is then applied to $A$ or $B$. For example, considering
$A = {\tt 0110\,0111}$ and $B = {\tt 0111\,1111}$,
$X = \lnot (A \oplus B) = {\tt 1110\,0111}$, setting each bit on the right side
of a 0 to 0 we get ${\tt 1110\,0000}$. We apply this mask to $A$ and get the
greatest common prefix of $A$ and $B$: ${\tt 011*\,****}$. This algorithm is
used by the first version of out implementation and is named {\em Bittree-naive}
in the experiments.

We optimized this algorithm by firstly computing all of the 65 or 33 different
masks (from {\tt 0x0} to {\tt0xf$\dots$f}) and storing them in an array. Then we
use a dichotomic search to find the mask corresponding to the greatest common
prefix: if $A$ and $B$ have a 32-bit common prefix, do they have a 48-bit common
prefix ? Otherwise, do they have a 16-bit common prefix ? And so on. This search
takes at most 6 (respectively 5) steps on 64-bit (respectively 32-bit) bittrees.
This algorithm is used by the final version of our implementation and is named
{\em Bittree-opti} in the experiments.


\subsection{Experiments}

These experiments justify the choice of bittrees over linked lists and binary
search trees. We implemented the classic {\em merge sort} algorithm, and added
extra allocations/deallocations to put each data structure to the test.
The program has been instrumented (see Section~\ref{section:instru}) to call our
monitoring functions (see Section~\ref{section:lib}). The execution time of the
instrumented program using each data structure has been measured (in
micro-seconds) and is plotted against the number of calls to a function $store$
(adding an element to the data structure). Figure~\ref{fig:experiments}
displays the results of the experiments.
The reference time is the execution time of the program without any
instrumentation. Bittree-naive and Bittree-opti are using two different versions
of the greatest common prefix computation (see Sub-section~\ref{section:gcpc}).



\begin{figure}[h]
  \hspace{-0.5in}
  \begin{tikzpicture}
    \begin{axis}[xlabel={Number of calls to $store$}, ylabel={Time ($\mu$s)},
	legend entries={List, Tree, Bittree-naive, Bittree-opti, Reference time}
      ]
      \addplot table {timeList.dat};
      \addplot table {timeTree.dat};
      \addplot table {timeBTn.dat};
      \addplot table {timeBTo.dat};
      \addplot table {timeNone.dat};
    \end{axis}
  \end{tikzpicture}
  \caption{Execution time plotted against the number of calls to $store$}
  \label{fig:experiments}
\end{figure}



The last but one measurement, the merge sort applied to 50,000 elements
(4,889,819 calls to $store$), has a reference time of 0.07s. It runs 32s with
the optimized bittree, 118s with the naive bittree, 7 hours 15 minutes with the
binary search trees and 19 hours 15 minutes with linked lists.
Our last measurement, the merge sort applied to 100,000 elements
(10,279,851 calls to $store$)  has a reference time of 0.19s. It runs 72s with
the optimized bittree, 252s with the naive bittree, but binary search trees and
linked lists exceed our 24-hour timeout.



\section{Monitoring functions}
\label{section:lib}


This section presents the functions used for adding/removing/retrieving
in formations about block descriptors into our data structure (bittree). These
functions will be automatically inserted in the source code by the
instrumentation step (see Section~\ref{section:instru}).


\subsection{Automatic allocation}

Automatic allocation functions keep track of all non-dynamically allocated
variables (but occupying memory space nevertheless), such as formal parameters,
local variables or global variables.

\begin{lstlisting}[caption=Automatic allocation functions]
void * _store_block
  (void * ptr, size_t size);
void _delete_block(void * ptr);
\end{lstlisting}


\subsection{Dynamic allocation}

These functions have to be used instead of those of the standard library 
({\em stdlib.h}).

\begin{lstlisting}[caption=Dynamic allocation functions]
void * _malloc(size_t size);
void * _realloc
  (void * ptr, size_t size);
void * _calloc
  (size_t nbr, size_t size);
void _free(void * ptr);
\end{lstlisting}


\subsection{Initialization}

These functions have to be used for each assignment, to update the
initialization status of a block descriptor. ${\tt \_initialize(ptr, size)}$
marks the $size$ first bytes starting from $ptr$ as initialized.
${\tt \_full\_init(ptr)}$ marks all the bytes of $ptr$ as initialized at once.
It is designed to avoid multiple calls to ${\tt \_initialize}$ whenever it is
possible and improves efficiency of the instrumented program.


\begin{lstlisting}[caption=Initialization functions]
void _initialize
  (void * ptr, size_t size);
void _full_init(void * ptr);
\end{lstlisting}


\subsection{Interrogation}

These functions are used to retrieve information about the block descriptors
and match the \textsc{E-Acsl} annotations we are trying to support:
${\tt \backslash valid}$, ${\tt \backslash base\_addr}$,
${\tt \backslash block\_length}$, ${\tt \backslash offset}$ and
${\tt \backslash initialized}$.

\begin{lstlisting}[caption=Interrogation functions]
int _valid
  (void * ptr, size_t size);
void * _base_addr(void * ptr);
size_t _block_length(void * ptr);
int _offset
  (void * ptr, size_t size);
int _initialized
  (void * ptr, size_t size);
\end{lstlisting}

${\tt \_valid(ptr, size)}$ returns 1 if it is safe to read/write $size$ bytes
starting from $ptr$, 0 otherwise. ${\tt \_base\_addr(ptr)}$ returns the base
address of the block containing $ptr$ if such a block exists, NULL otherwise.
${\tt \_block\_length(ptr)}$ returns the size (in bytes) of the block containing
$ptr$ if such a block exists, 0 otherwise. ${\tt \_offset(ptr)}$ returns the
offset between $ptr$ and the base address of the block containing $ptr$ if such
a block exists, -1 otherwise. ${\tt \_initialized(ptr, size)}$ returns 1 if the
$size$ first bytes starting from $ptr$ are initialized, 0 otherwise.


\section{Instrumentation}
\label{section:instru}

This section presents the instrumentation performed by the \textsc{E-Acsl}
plug-in to monitor the memory used by a program. The generated instrumented code
uses the previously defined functions (see previous section).

For each global variable, calls to ${\tt \_store\_block}$ and
${\tt \_full\_init}$ are inserted at the beginning of the $main$ function, and a
call to ${\tt \_delete\_block}$ at the end (see Figure~\ref{fig:globals}). For
each formal parameter and local variable, a call to ${\tt \_store\_block}$ is
inserted at the beginning of their scope block and a call to
${\tt \_delete\_block}$ at the end of their scope block (see
Figure~\ref{fig:formals} and Figure~\ref{fig:locals}). Calls to
${\tt \_full\_init}$ and ${\tt \_initialize}$ are inserted on assignments (see
Figure~\ref{fig:init}), and \textsc{E-Acsl} annotations are translated to the
corresponding functions which result is tested by an assertion (see
Figure~\ref{fig:annot}).



\begin{figure}[h]
    \begin{lstlisting}
int * p;
p = §\highl{\_malloc}§(32);
§\highl{\_free}§(p);
    \end{lstlisting}
  \caption{Dynamic allocation instrumentation}
  \label{fig:alloc-dyn}
\end{figure}


\begin{figure}[h]
    \begin{lstlisting}
{
  int* p;
  §\highl{\_store\_block(\&p, sizeof(int*));}§
  ...
  §\highl{\_delete\_block(\&p);}§
}
    \end{lstlisting}
  \caption{Local variable instrumentation}
  \label{fig:locals}
\end{figure}


\begin{figure}[h]
    \begin{lstlisting}
int g;

int main() {
  §\highl{\_store\_block(\&g, sizeof(int));}§
  §\highl{\_full\_init(\&g);}§
  ...
  §\highl{\_delete\_block(\&g);}§
  ...
  return 0;
}
    \end{lstlisting}
  \caption{Global variable instrumentation}
  \label{fig:globals}
\end{figure}


\begin{figure}[h]
    \begin{lstlisting}
void f (int i) {
  §\highl{\_store\_block(\&i);}§
  §\highl{\_full\_init(\&i);}§
  ...
  §\highl{\_delete\_block(\&i);}§
}
    \end{lstlisting}
  \caption{Formal parameter instrumentation}
  \label{fig:formals}
\end{figure}


\begin{figure}[h]
    \begin{lstlisting}
int i;
i = 4;
§\highl{\_full\_init(\&i);}§

int t [10];
t[2] = 4;
§\highl{\_initialize((t+2), sizeof(int));}§
    \end{lstlisting}
  \caption{Assignment instrumentation}
  \label{fig:init}
\end{figure}


\begin{figure}[h]
    \begin{lstlisting}
int p[12];
//@ assert \valid (p+2);
§\highl{assert(\_valid((p+2), sizeof(int)));}§
    \end{lstlisting}
  \caption{Annotation instrumentation}
  \label{fig:annot}
\end{figure}


\section{Conclusion}


We implemented an efficient data structure (bittree) (see
Section~\ref{section:bittree}) to store the block descriptors and functions (see
Section~\ref{section:lib}) relying on this structure to perform memory
monitoring.
We also defined and implemented into the \textsc{E-Acsl} plug-in the
instrumentation (see Section~\ref{section:instru}) to perform on a C source code
to monitor the memory. This allows the runtime checking of the following
\textsc{E-Acsl} annotations: ${\tt \backslash base\_addr}$,
${\tt \backslash block\_length}$, ${\tt \backslash offset}$,
${\tt \backslash valid}$ and ${\tt \backslash initialized}$.


\begin{scriptsize}
\bibliographystyle{abbrv}
\bibliography{biblio}
\end{scriptsize}





\end{document}

